\documentclass[11pt,a4paper]{article}
\usepackage[margin=1in]{geometry}
\usepackage[T1]{fontenc}
\usepackage[utf8]{inputenc}
\usepackage{lmodern}
\usepackage{setspace}
\usepackage{parskip}
\usepackage{hyperref}
\usepackage{longtable}
\usepackage{booktabs}
\usepackage{array}
\usepackage{xcolor}
\usepackage{listings}
\usepackage{enumitem}
\usepackage{titlesec}
\usepackage{fancyhdr}
\usepackage{graphicx}

\hypersetup{
  colorlinks=true,
  linkcolor=blue!60!black,
  urlcolor=blue!60!black,
  pdftitle={University Accommodation Office System - CS266 Presentation Report},
  pdfauthor={Team 24bcs001-24bcs007}
}

\setstretch{1.1}
\setlist[itemize]{leftmargin=1.4em}
\setlist[enumerate]{leftmargin=1.6em}

\titleformat{\section}{\large\bfseries}{\thesection}{0.6em}{}
\titleformat{\subsection}{\normalsize\bfseries}{\thesubsection}{0.6em}{}

\pagestyle{fancy}
\fancyhf{}
\fancyhead[L]{CS266 DBMS Project Report}
\fancyhead[R]{University Accommodation Office System}
\fancyfoot[C]{\thepage}
\setlength{\headheight}{14pt}

\lstdefinestyle{codeblock}{
  basicstyle=\ttfamily\small,
  breaklines=true,
  frame=single,
  rulecolor=\color{black!25},
  backgroundcolor=\color{black!2},
  keywordstyle=\color{blue!65!black},
  commentstyle=\color{green!40!black},
  stringstyle=\color{orange!60!black},
  showstringspaces=false,
  tabsize=2
}

\lstset{style=codeblock}

\begin{document}

\begin{titlepage}
\centering
\vspace*{1.2cm}

{\LARGE\bfseries CS266 - Database Management Systems\par}
\vspace{0.6cm}
{\Large\bfseries University Accommodation Office Management System\par}
\vspace{0.4cm}
{\large Presentation and Technical Report (LaTeX)\par}

\vspace{1.0cm}

\begin{tabular}{ll}
\textbf{Project Type:} & Full-stack DBMS Application \\
\textbf{Frontend:} & React + TypeScript + Vite \\
\textbf{Backend:} & FastAPI + SQLAlchemy Core \\
\textbf{Database:} & MySQL 8+ \\
\end{tabular}

\vspace{1.2cm}

{\large\bfseries Team Roll Numbers\par}
\vspace{0.4cm}
\begin{tabular}{c}
24bcs001 \\
24bcs002 \\
24bcs003 \\
24bcs004 \\
24bcs005 \\
24bcs006 \\
24bcs007 \\
\end{tabular}

\vspace{1.4cm}

{\large Date: \today\par}

\vfill
\end{titlepage}

\tableofcontents
\newpage

\section{Executive Summary}
This project is a complete University Accommodation Office Management System designed as a DBMS-centric full-stack application. It solves a practical institutional problem: tracking students, staff, housing inventory, leases, invoices, compliance inspections, and emergency contacts in one consistent relational system.

The implementation is split into three layers:
\begin{itemize}
  \item \textbf{Database layer (MySQL):} normalized schema, constraints, triggers, indexes, and report views.
  \item \textbf{Service layer (FastAPI + SQLAlchemy):} secure APIs for authentication, CRUD, and reporting.
  \item \textbf{Presentation layer (React + TypeScript):} login, entity studio, dashboards, and report execution UI.
\end{itemize}

This report explains the project from scratch to deployment, including every process flow, all API operations, DB connectivity, transaction handling, and how ACID properties are realized in practice.

\section{Problem Statement and Objectives}
\subsection{Problem Context}
Accommodation offices often suffer from fragmented data and manual workflows:
\begin{itemize}
  \item Student allocation data in spreadsheets.
  \item Lease and invoice records in disconnected systems.
  \item Slow or inconsistent reporting.
  \item Weak integrity checks that allow invalid data.
\end{itemize}

\subsection{Project Objectives}
\begin{enumerate}
  \item Build a normalized relational schema for accommodation operations.
  \item Provide secure, role-based APIs for operational workflows.
  \item Support complete CRUD across all key entities.
  \item Implement DBMS-style report queries via SQL views.
  \item Provide a clean, modern frontend for daily usage.
  \item Ensure data consistency through constraints, triggers, and transaction handling.
\end{enumerate}

\section{Complete Tech Stack}
\subsection{Frontend}
\begin{itemize}
  \item React 18.3.1
  \item TypeScript 5.7.2
  \item Vite 6.2.1
  \item React Router DOM 6.30.0
  \item Tailwind CSS 3.4.17
  \item Framer Motion 12.9.4
  \item Lucide React 0.469.0
\end{itemize}

\subsection{Backend}
\begin{itemize}
  \item FastAPI 0.115.12
  \item Uvicorn 0.34.0
  \item SQLAlchemy 2.0.40 (Core style + reflection)
  \item PyMySQL 1.1.1
  \item python-dotenv 1.0.1
  \item python-jose 3.4.0
  \item passlib 1.7.4
\end{itemize}

\subsection{Database and Testing}
\begin{itemize}
  \item MySQL 8+ (or compatible managed MySQL)
  \item pytest 8.3.5
  \item httpx 0.28.1
\end{itemize}

\section{Repository Structure and Responsibilities}
\begin{lstlisting}[language={},caption={Project layout at high level}]
DBMS_courseproject/
  backend/
    app/
      main.py       # FastAPI routes: auth, reports, generic CRUD
      auth.py       # JWT auth, role checks, seeded users
      config.py     # environment loading and settings parsing
      database.py   # engine, session, metadata reflection
    schema.sql      # tables, FKs, triggers, indexes, views
    requirements.txt
    tests/
  frontend/
    src/
      lib/api.ts          # fetch wrapper + token injection
      providers/AuthProvider.tsx
      config/entities.ts  # UI metadata for all entities
      config/reports.ts   # report catalog and parameters
      pages/
    vite.config.ts
  seed.sql
  docs/
\end{lstlisting}

\section{System Architecture from Start to End}
\subsection{High-Level Flow}
\begin{enumerate}
  \item User interacts with frontend page (login, CRUD, reports).
  \item Frontend calls API helper in \texttt{frontend/src/lib/api.ts}.
  \item Request reaches FastAPI route in \texttt{backend/app/main.py}.
  \item Dependencies validate token and role (if required).
  \item Backend executes SQL via SQLAlchemy session.
  \item MySQL returns result or constraint error.
  \item Backend serializes response into JSON.
  \item Frontend renders result in UI.
\end{enumerate}

\subsection{Development Proxy and Service Boundaries}
Vite proxy routes local requests to backend:
\begin{lstlisting}[language={},caption={Vite dev proxy}]
proxy: {
  "/api": "http://localhost:8000",
  "/health": "http://localhost:8000",
  "/auth": "http://localhost:8000"
}
\end{lstlisting}

\section{Database Design (Relational Core)}
\subsection{Primary Entities}
The core schema contains 10 transactional tables:
\begin{enumerate}
  \item staff
  \item courses
  \item students
  \item halls
  \item apartments
  \item rooms
  \item leases
  \item invoices
  \item next\_of\_kin
  \item inspections
\end{enumerate}

\subsection{ER Diagram Placeholder}
\begin{figure}[h]
\centering
\fbox{%
  \parbox[c][0.28\textheight][c]{0.90\textwidth}{%
    \centering
    \textbf{ER Diagram Placeholder}\\[0.6em]
    Add your ER diagram image file at: \texttt{docs/figures/erd.png}\\
    Replace this box with:\\
    \texttt{\textbackslash includegraphics[width=0.95\textbackslash textwidth]\{docs/figures/erd.png\}}
  }%
}
\caption{Entity Relationship Diagram (to be inserted)}
\label{fig:erd-placeholder}
\end{figure}

\subsection{Important Relational Rules}
\begin{itemize}
  \item \textbf{students.course\_number} references courses with \texttt{ON DELETE SET NULL}.
  \item \textbf{students.advisor\_staff\_id} references staff with \texttt{ON DELETE SET NULL}.
  \item \textbf{leases.banner\_id} references students with \texttt{ON DELETE CASCADE}.
  \item \textbf{invoices.lease\_id} references leases with \texttt{ON DELETE CASCADE}.
  \item \textbf{next\_of\_kin.banner\_id} is UNIQUE, enforcing max one kin row per student.
\end{itemize}

\subsection{Data Integrity Controls}
\begin{itemize}
  \item Check constraints:
  \begin{itemize}
    \item apartment bedrooms in \{3,4,5\}
    \item room monthly\_rent \(\ge 0\)
    \item invoice amount\_due \(\ge 0\)
  \end{itemize}
  \item Trigger-based logic in \texttt{rooms}:
  \begin{itemize}
    \item exactly one of hall\_id or apartment\_id must be non-null
    \item enforced on both INSERT and UPDATE
  \end{itemize}
\end{itemize}

\subsection{Report Views (DBMS-Oriented)}
The schema defines 14 views for report routes:
\begin{itemize}
  \item v\_hall\_managers, v\_student\_leases, v\_summer\_leases
  \item v\_student\_rent\_paid, v\_unpaid\_invoices
  \item v\_unsatisfactory\_inspections, v\_hall\_student\_rooms
  \item v\_waiting\_list, v\_student\_category\_counts
  \item v\_students\_without\_kin, v\_student\_advisers
  \item v\_hall\_rent\_stats, v\_hall\_place\_counts, v\_senior\_staff
\end{itemize}

\section{How Backend is Connected to DB}
\subsection{Environment Loading}
\begin{lstlisting}[language=Python,caption={Environment initialization in config.py}]
BACKEND_ROOT = Path(__file__).resolve().parents[1]
load_dotenv(BACKEND_ROOT / ".env")
load_dotenv(override=False)
\end{lstlisting}

\subsection{DB Settings Used}
\begin{itemize}
  \item \texttt{DB\_HOST}
  \item \texttt{DB\_PORT}
  \item \texttt{DB\_USER}
  \item \texttt{DB\_PASSWORD}
  \item \texttt{DB\_NAME}
  \item \texttt{DB\_TLS\_MODE} (false, preferred, skip-verify)
  \item \texttt{DB\_CA\_CERT\_PATH}
\end{itemize}

\subsection{SQLAlchemy + PyMySQL Connection URL}
\begin{lstlisting}[language=Python,caption={Connection URL construction}]
return (
    f"mysql+pymysql://{user}:{password}@{settings.db_host}:{settings.db_port}/"
    f"{settings.db_name}?charset=utf8mb4"
)
\end{lstlisting}

\subsection{Engine and Session Lifecycle}
\begin{itemize}
  \item Lazy singleton engine creation in \texttt{get\_engine()}.
  \item Connection resilience via \texttt{pool\_pre\_ping=True} and \texttt{pool\_recycle=300}.
  \item Per-request session from \texttt{sessionmaker}, then session closed in \texttt{finally}.
\end{itemize}

\subsection{Metadata Reflection Strategy}
Instead of hardcoding ORM models per table, backend reflects metadata at runtime and fetches tables by name. This enables generic CRUD for all entities with one route pattern.

\section{Authentication and Authorization Internals}
\subsection{Auth Model}
\begin{itemize}
  \item OAuth2 bearer token in request header.
  \item JWT payload contains \texttt{sub}, \texttt{role}, \texttt{exp}.
  \item Password verification with passlib \texttt{pbkdf2\_sha256}.
\end{itemize}

\subsection{Seeded Demo Accounts}
\begin{itemize}
  \item admin / Admin@123 (role: admin)
  \item manager / Manager@123 (role: manager)
  \item viewer / Viewer@123 (role: viewer)
\end{itemize}

\subsection{Role Guards}
\begin{itemize}
  \item READ roles: admin, manager, viewer
  \item WRITE roles: admin, manager
  \item DELETE role: admin only
\end{itemize}

\section{Frontend Process Flow from Scratch}
\subsection{Router and Guards}
\begin{itemize}
  \item \texttt{/login} is wrapped in \texttt{LoginOnly}.
  \item Main app routes are wrapped in \texttt{RequireAuth}.
  \item Unauthorized users are redirected to login.
\end{itemize}

\subsection{API Wrapper Behavior}
\begin{lstlisting}[language={},caption={Auth token injection in api.ts}]
if (!skipAuth && authToken && !headers.has("Authorization")) {
  headers.set("Authorization", `Bearer ${authToken}`);
}
\end{lstlisting}

\subsection{Session Persistence}
AuthProvider stores token and user in localStorage and refreshes profile from \texttt{/auth/me} on app bootstrap.

\section{Complete API Inventory and Operations}
\subsection{Core Endpoints}
\begin{longtable}{>{\raggedright\arraybackslash}p{0.12\textwidth} >{\raggedright\arraybackslash}p{0.36\textwidth} >{\raggedright\arraybackslash}p{0.16\textwidth} >{\raggedright\arraybackslash}p{0.28\textwidth}}
\toprule
\textbf{Method} & \textbf{Endpoint} & \textbf{Access} & \textbf{Purpose} \\
\midrule
\endhead
GET & /health & Public & Service health probe \\
POST & /auth/login & Public & Issue JWT token \\
GET & /auth/me & Any authenticated role & Current user profile \\
GET & /api/\{entity\} & read roles & List records \\
GET & /api/\{entity\}/\{id\} & read roles & Get single record \\
POST & /api/\{entity\} & write roles & Create record \\
PUT & /api/\{entity\}/\{id\} & write roles & Update record \\
DELETE & /api/\{entity\}/\{id\} & admin & Delete record \\
\bottomrule
\end{longtable}

Supported entities in generic CRUD:
\begin{itemize}
  \item staff, courses, students, halls, apartments
  \item rooms, leases, invoices, next-of-kin, inspections
\end{itemize}

\subsection{Report Endpoints (a to n)}
\begin{longtable}{>{\raggedright\arraybackslash}p{0.05\textwidth} >{\raggedright\arraybackslash}p{0.45\textwidth} >{\raggedright\arraybackslash}p{0.42\textwidth}}
\toprule
\textbf{ID} & \textbf{Endpoint} & \textbf{SQL View} \\
\midrule
\endhead
a & /api/reports/hall-managers & v\_hall\_managers \\
b & /api/reports/student-leases & v\_student\_leases \\
c & /api/reports/summer-leases & v\_summer\_leases \\
d & /api/reports/student-rent-paid/\{banner\_id\} & v\_student\_rent\_paid \\
e & /api/reports/unpaid-invoices?due\_before=... & v\_unpaid\_invoices \\
f & /api/reports/unsatisfactory-inspections & v\_unsatisfactory\_inspections \\
g & /api/reports/hall-student-rooms/\{hall\_id\} & v\_hall\_student\_rooms \\
h & /api/reports/waiting-list & v\_waiting\_list \\
i & /api/reports/student-category-counts & v\_student\_category\_counts \\
j & /api/reports/students-without-kin & v\_students\_without\_kin \\
k & /api/reports/student-adviser/\{banner\_id\} & v\_student\_advisers \\
l & /api/reports/rent-stats & v\_hall\_rent\_stats \\
m & /api/reports/hall-place-counts & v\_hall\_place\_counts \\
n & /api/reports/senior-staff & v\_senior\_staff \\
\bottomrule
\end{longtable}

\section{API Calling Examples (Complete Operational Coverage)}
\subsection{Step 1: Login and Capture Token}
\begin{lstlisting}[language={},caption={Authenticate and obtain bearer token}]
curl -X POST http://localhost:8000/auth/login \
  -H "Content-Type: application/json" \
  -d '{"username":"admin","password":"Admin@123"}'
\end{lstlisting}

\subsection{Step 2: Read Current User Profile}
\begin{lstlisting}[language={}]
curl http://localhost:8000/auth/me \
  -H "Authorization: Bearer <TOKEN>"
\end{lstlisting}

\subsection{Step 3: CRUD Operations Example (students)}
\begin{lstlisting}[language={},caption={Create student}]
curl -X POST http://localhost:8000/api/students \
  -H "Authorization: Bearer <TOKEN>" \
  -H "Content-Type: application/json" \
  -d '{
    "banner_id":"B00999999",
    "first_name":"Aanya",
    "last_name":"Sharma",
    "street":"1 Main St",
    "city":"London",
    "postcode":"SW1A 1AA",
    "dob":"2003-06-21",
    "gender":"Female",
    "category":"Postgraduate",
    "nationality":"Indian",
    "status":"Waiting",
    "major":"Computer Science"
  }'
\end{lstlisting}

\begin{lstlisting}[language={},caption={List students}]
curl http://localhost:8000/api/students \
  -H "Authorization: Bearer <TOKEN>"
\end{lstlisting}

\begin{lstlisting}[language={},caption={Get one student}]
curl http://localhost:8000/api/students/B00999999 \
  -H "Authorization: Bearer <TOKEN>"
\end{lstlisting}

\begin{lstlisting}[language={},caption={Update student}]
curl -X PUT http://localhost:8000/api/students/B00999999 \
  -H "Authorization: Bearer <TOKEN>" \
  -H "Content-Type: application/json" \
  -d '{"status":"Placed"}'
\end{lstlisting}

\begin{lstlisting}[language={},caption={Delete student (admin only)}]
curl -X DELETE http://localhost:8000/api/students/B00999999 \
  -H "Authorization: Bearer <TOKEN>"
\end{lstlisting}

\subsection{Step 4: Report API Examples}
\begin{lstlisting}[language={},caption={Report e: unpaid invoices up to date}]
curl "http://localhost:8000/api/reports/unpaid-invoices?due_before=2026-12-31" \
  -H "Authorization: Bearer <TOKEN>"
\end{lstlisting}

\begin{lstlisting}[language={},caption={Report g: hall student rooms}]
curl http://localhost:8000/api/reports/hall-student-rooms/1 \
  -H "Authorization: Bearer <TOKEN>"
\end{lstlisting}

\begin{lstlisting}[language={},caption={Report n: senior staff}]
curl http://localhost:8000/api/reports/senior-staff \
  -H "Authorization: Bearer <TOKEN>"
\end{lstlisting}

\section{Detailed Internal Process Flows}
\subsection{Flow A: Login Process}
\begin{enumerate}
  \item User submits credentials from \texttt{LoginPage}.
  \item Frontend calls \texttt{apiPost("/auth/login", ...)} with \texttt{skipAuth=true}.
  \item FastAPI route \texttt{/auth/login} calls \texttt{authenticate\_user()}.
  \item Password is verified with passlib hash.
  \item On success, JWT is generated by \texttt{create\_access\_token()}.
  \item Frontend persists token + profile and sets API auth token in memory.
\end{enumerate}

\subsection{Flow B: Generic Create Operation}
\begin{enumerate}
  \item Request hits \texttt{POST /api/\{entity\}}.
  \item Entity key is validated against \texttt{TABLE\_CONFIG}.
  \item SQLAlchemy reflected table is fetched with \texttt{get\_table()}.
  \item Payload is sanitized and type-normalized by column type.
  \item SQL \texttt{INSERT} is executed.
  \item \texttt{\_safe\_commit()} commits transaction or rolls back on error.
  \item Created row is re-fetched and returned.
\end{enumerate}

\subsection{Flow B1: Generic List Route for Each Entity (\texttt{GET /api/\{entity\}})}
Exact backend implementation:
\begin{lstlisting}[language=Python,caption={Generic list route used for every entity}]
@app.get("/api/{entity}")
def list_records(
    entity: str,
    _: AuthUser = Depends(READ_ROLES),
    session: Session = Depends(get_session),
) -> Any:
    config = _entity_config(entity)
    table = get_table(config["table"])

    query = select(table)
    pk_name = config["pk"]
    if pk_name in table.c:
        query = query.order_by(table.c[pk_name])

    rows = session.execute(query).mappings().all()
    return _encode([dict(row) for row in rows])
\end{lstlisting}

Entity-specific behavior is derived from \texttt{TABLE\_CONFIG}:
\begin{longtable}{>{\raggedright\arraybackslash}p{0.15\textwidth} >{\raggedright\arraybackslash}p{0.16\textwidth} >{\raggedright\arraybackslash}p{0.14\textwidth} >{\raggedright\arraybackslash}p{0.12\textwidth} >{\raggedright\arraybackslash}p{0.33\textwidth}}
\toprule
\textbf{Entity URL Key} & \textbf{DB Table} & \textbf{PK Used for Sort} & \textbf{ID Type} & \textbf{Runtime Result in List Route} \\
\midrule
\endhead
staff & staff & staff\_id & int & \texttt{SELECT * FROM staff ORDER BY staff\_id} \\
courses & courses & course\_number & str & \texttt{SELECT * FROM courses ORDER BY course\_number} \\
students & students & banner\_id & str & \texttt{SELECT * FROM students ORDER BY banner\_id} \\
halls & halls & hall\_id & int & \texttt{SELECT * FROM halls ORDER BY hall\_id} \\
apartments & apartments & apartment\_id & int & \texttt{SELECT * FROM apartments ORDER BY apartment\_id} \\
rooms & rooms & place\_number & int & \texttt{SELECT * FROM rooms ORDER BY place\_number} \\
leases & leases & lease\_id & int & \texttt{SELECT * FROM leases ORDER BY lease\_id} \\
invoices & invoices & invoice\_id & int & \texttt{SELECT * FROM invoices ORDER BY invoice\_id} \\
next-of-kin & next\_of\_kin & kin\_id & int & \texttt{SELECT * FROM next\_of\_kin ORDER BY kin\_id} \\
inspections & inspections & inspection\_id & int & \texttt{SELECT * FROM inspections ORDER BY inspection\_id} \\
\bottomrule
\end{longtable}

Step-by-step runtime behavior for each request:
\begin{enumerate}
  \item \texttt{entity} path value is validated by \texttt{\_entity\_config(entity)}.
  \item Table name and PK are resolved from \texttt{TABLE\_CONFIG}.
  \item Reflected table object is loaded via \texttt{get\_table(config["table"])}.
  \item \texttt{select(table)} query object is created.
  \item PK-based deterministic ordering is appended: \texttt{ORDER BY <pk>}.
  \item SQL is executed with \texttt{session.execute(query)}.
  \item Rows are materialized as mapping dictionaries.
  \item Payload is JSON-encoded via \texttt{\_encode(...)} and returned to frontend.
\end{enumerate}

\subsection{Flow C: Report Query Execution}
\begin{enumerate}
  \item Frontend chooses report metadata from \texttt{reports.ts}.
  \item Endpoint is built (path and query params).
  \item FastAPI report route runs SQL text against a view.
  \item \texttt{session.execute(text(sql), params)} returns row mappings.
  \item Rows are converted to dictionaries and encoded as JSON.
\end{enumerate}

\section{How Type Normalization and Validation Work}
Before writes, backend coerces values according to SQL type:
\begin{itemize}
  \item Date columns: strict \texttt{YYYY-MM-DD}
  \item DateTime columns: ISO datetime parsing
  \item Boolean columns: true/false, 1/0, yes/no parsing
  \item Integer and numeric fields: explicit casting
\end{itemize}

If parsing fails, API returns 400 with field-specific message.

\section{Error Handling Strategy}
\subsection{Constraint and Integrity Mapping}
\texttt{\_map\_integrity\_error} maps common MySQL errors:
\begin{itemize}
  \item 1062: duplicate key \(\rightarrow\) 409 conflict
  \item 1452: invalid foreign key reference \(\rightarrow\) 400 bad request
  \item 3819 / 1644 / 1406: validation or trigger issues \(\rightarrow\) 400 bad request
\end{itemize}

\subsection{Transaction Safety}
\texttt{\_safe\_commit()} guarantees:
\begin{itemize}
  \item commit on success
  \item rollback on any SQLAlchemy error
  \item mapped HTTP errors returned to client
\end{itemize}

\section{ACID Properties: Theory and Implementation in This Project}
\subsection{Important Clarification}
ACID is guaranteed by the database engine (MySQL/InnoDB), not by SQLAlchemy alone. SQLAlchemy provides transaction management APIs that correctly use DB transactions.

\subsection{A - Atomicity}
\begin{itemize}
  \item Each write request is transactional.
  \item If any step fails, rollback happens, so partial writes are avoided.
\end{itemize}

\subsection{C - Consistency}
\begin{itemize}
  \item Enforced by schema constraints, FKs, checks, and triggers.
  \item Invalid states (e.g., room with both hall and apartment) are blocked.
\end{itemize}

\subsection{I - Isolation}
\begin{itemize}
  \item MySQL transaction isolation is applied per connection/session.
  \item Concurrent transactions are isolated according to DB isolation level.
\end{itemize}

\subsection{D - Durability}
\begin{itemize}
  \item Once commit succeeds, data persists in MySQL storage engine.
  \item Server restarts do not lose committed data.
\end{itemize}

\subsection{Practical Note on SQLAlchemy}
SQLAlchemy provides:
\begin{itemize}
  \item session boundaries
  \item commit and rollback orchestration
  \item connection pooling and retries
\end{itemize}
But data durability and strict ACID guarantees are still fundamentally DB responsibilities.

\section{Code-Level Examples from Real Project}
\subsection{FastAPI Route Example}
\begin{lstlisting}[language=Python,caption={Report route structure}]
@app.get("/api/reports/hall-managers")
def report_hall_managers(
    _: AuthUser = Depends(READ_ROLES),
    session: Session = Depends(get_session),
) -> Any:
    sql = """
        SELECT hall_id, hall_name, manager_name, manager_phone
        FROM v_hall_managers
        ORDER BY hall_id
    """
    return _encode(_run_report(session, sql))
\end{lstlisting}

\subsection{Generic CRUD Example}
\begin{lstlisting}[language=Python,caption={List route used for all entities}]
@app.get("/api/{entity}")
def list_records(entity: str, _: AuthUser = Depends(READ_ROLES), session: Session = Depends(get_session)) -> Any:
    config = _entity_config(entity)
    table = get_table(config["table"])
    query = select(table)
    rows = session.execute(query).mappings().all()
    return _encode([dict(row) for row in rows])
\end{lstlisting}

\subsection{Frontend API Wrapper Example}
\begin{lstlisting}[language={},caption={Single API helper used by all pages}]
export async function apiRequest<T>(path: string, init: ApiRequestOptions = {}): Promise<T> {
  const { skipAuth = false, ...requestInit } = init;
  const headers = new Headers(requestInit.headers ?? {});

  if (!skipAuth && authToken && !headers.has("Authorization")) {
    headers.set("Authorization", `Bearer ${authToken}`);
  }

  const response = await fetch(buildUrl(path), { ...requestInit, headers });
  const payload = await parseResponse(response);
  if (!response.ok) {
    throw new ApiError(response.status, asErrorMessage(payload, `Request failed with status ${response.status}`));
  }
  return payload as T;
}
\end{lstlisting}

\section{Report Layer and DBMS Demonstration Value}
This project is strong for DBMS presentation because report logic is encoded in SQL views and then consumed by API endpoints. This allows demonstration in both places:
\begin{itemize}
  \item Workbench: direct \texttt{SELECT * FROM v\_...}
  \item Application: authenticated \texttt{/api/reports/...}
\end{itemize}

\section{Deep Execution Pipeline: From \texttt{app.get(...)} to \texttt{session.execute(...)}}
\subsection{ASGI-to-Route Lifecycle (Exact Runtime Sequence)}
For every request handled by a FastAPI route (for example \texttt{@app.get("/api/reports/hall-managers")}), the runtime path is:
\begin{enumerate}
  \item Uvicorn receives HTTP request and forwards it to the FastAPI ASGI app.
  \item Starlette/FastAPI router matches path and method to a concrete handler.
  \item FastAPI builds dependency graph for that handler.
  \item Security dependency extracts bearer token (if endpoint is protected).
  \item JWT is decoded and role is validated by \texttt{require\_roles(...)}.
  \item DB dependency \texttt{get\_session()} yields a SQLAlchemy \texttt{Session}.
  \item Route function executes Python logic and builds SQL expression/text.
  \item \texttt{session.execute(...)} sends SQL through engine to PyMySQL driver.
  \item MySQL executes query and returns rows/metadata.
  \item Result is materialized with \texttt{.mappings().all()} or \texttt{.first()}.
  \item Backend encodes payload through \texttt{jsonable\_encoder} and returns HTTP response.
  \item Session dependency finalizer closes DB session.
\end{enumerate}

\subsection{How \texttt{app.get(...)} is Called Internally}
When a route is declared with \texttt{@app.get("/api/...")}, FastAPI registers metadata for method/path/handler. At runtime, the decorator itself is not called again; instead, routing metadata is used to dispatch incoming requests to the associated function.

\subsection{CRUD Query Object Construction (SQLAlchemy Core)}
In generic CRUD routes, backend builds SQL objects before execution. Example pattern:
\begin{lstlisting}[language=Python,caption={SQLAlchemy query object pipeline in CRUD}]
config = _entity_config(entity)
table = get_table(config["table"])

query = select(table)
pk_name = config["pk"]
if pk_name in table.c:
    query = query.order_by(table.c[pk_name])

rows = session.execute(query).mappings().all()
\end{lstlisting}

Here, \texttt{query} is a SQLAlchemy expression object. It is compiled by the selected SQL dialect and then passed to the DBAPI driver at execution time.

\subsection{Report Query Construction (Textual SQL + Bind Params)}
In report routes, backend keeps SQL explicit and uses bind parameters:
\begin{lstlisting}[language=Python,caption={Report execution helper path}]
def _run_report(session: Session, sql: str, params: dict[str, Any] | None = None):
    rows = session.execute(text(sql), params or {}).mappings().all()
    return [dict(row) for row in rows]
\end{lstlisting}

This path preserves query readability for DBMS demos while still using SQLAlchemy session management and parameter binding.

\subsection{Write Operations and Transactions}
For \texttt{POST/PUT/DELETE}, the execution pattern is:
\begin{enumerate}
  \item Build \texttt{insert/update/delete} object.
  \item Execute with \texttt{session.execute(...)}.
  \item Call \texttt{\_safe\_commit()}.
  \item On integrity/SQL errors: rollback + mapped HTTP exception.
\end{enumerate}

\section{All Views Explained in Detail (Complete Report Layer)}
\begingroup
\scriptsize
\setlength{\tabcolsep}{3pt}
\renewcommand{\arraystretch}{1.08}
\begin{longtable}{>{\raggedright\arraybackslash}p{0.18\textwidth} >{\raggedright\arraybackslash}p{0.22\textwidth} >{\raggedright\arraybackslash}p{0.29\textwidth} >{\raggedright\arraybackslash}p{0.25\textwidth}}
\toprule
\textbf{View} & \textbf{API Route} & \textbf{Columns Returned} & \textbf{Definition Logic / Source Tables} \\
\midrule
\endhead
v\_hall\_managers & /api/reports/hall-managers & hall\_id, hall\_name, manager\_name, manager\_phone & \texttt{halls} JOIN \texttt{staff} on manager\_staff\_id; no filter \\

v\_student\_leases & /api/reports/student-leases & lease\_id, banner\_id, student\_name, duration\_semesters, includes\_summer\_semester, date\_enter, date\_leave, place\_number, room\_number, monthly\_rent, residence\_type, residence\_name, residence\_address & \texttt{leases} + \texttt{students} + \texttt{rooms} + optional \texttt{halls/apartments} joins \\

v\_summer\_leases & /api/reports/summer-leases & Same as \texttt{v\_student\_leases} & Derived from \texttt{v\_student\_leases}; filter \texttt{includes\_summer\_semester = TRUE} \\

v\_student\_rent\_paid & /api/reports/student-rent-paid/\{banner\_id\} & banner\_id, student\_name, total\_paid & LEFT JOIN chain from \texttt{students -> leases -> invoices}; conditional SUM on paid invoices \\

v\_unpaid\_invoices & /api/reports/unpaid-invoices & invoice\_id, lease\_id, banner\_id, student\_name, semester, amount\_due, due\_date, place\_number, room\_number, residence\_type, residence\_address & \texttt{invoices} JOIN \texttt{leases/students/rooms} + optional residence joins; filter \texttt{date\_paid IS NULL} \\

v\_unsatisfactory\newline inspections & /api/reports/\newline unsatisfactory-inspections & inspection\_id, inspection\_date, staff\_name, apartment\_id, comments & \texttt{inspections} JOIN \texttt{staff}; filter \texttt{is\_satisfactory = FALSE} \\

v\_hall\newline student\newline rooms & /api/reports/\newline hall-student-rooms/\{hall\_id\} & hall\_id, banner\_id, student\_name, hall\_name, room\_number, place\_number & \texttt{leases} JOIN \texttt{students/rooms/halls}; hall-scoped projection \\

v\_waiting\_list & /api/reports/waiting-list & banner\_id, first\_name, last\_name, street, city, postcode, mobile\_phone, email, dob, gender, category, nationality, special\_needs, comments, status, major, minor, course\_number, advisor\_staff\_id & Direct projection from \texttt{students}; filter \texttt{status='Waiting'} \\

v\_student\newline category\newline counts & /api/reports/\newline student-category-counts & category, student\_count & Group-by aggregate on \texttt{students.category} \\

v\_students\newline without\newline kin & /api/reports/\newline students-without-kin & banner\_id, student\_name & \texttt{students} LEFT JOIN \texttt{next\_of\_kin}; filter rows with null kin\_id \\

v\_student\_advisers & /api/reports/student-adviser/\{banner\_id\} & banner\_id, student\_name, adviser\_name, adviser\_phone & \texttt{students} LEFT JOIN adviser alias of \texttt{staff} \\

v\_hall\_rent\_stats & /api/reports/rent-stats & min\_rent, max\_rent, avg\_rent & Aggregate over \texttt{rooms} where \texttt{hall\_id IS NOT NULL} \\

v\_hall\_place\_counts & /api/reports/hall-place-counts & hall\_id, hall\_name, total\_places & \texttt{halls} LEFT JOIN \texttt{rooms}; grouped by hall \\

v\_senior\_staff & /api/reports/senior-staff & staff\_id, staff\_name, age, location & \texttt{staff} with \texttt{TIMESTAMPDIFF(YEAR, dob, CURDATE()) > 60} \\
\bottomrule
\end{longtable}
\endgroup

\section{Complete Per-Table Technical Specification (All 10 Tables)}
\subsection{\texttt{staff}}
\begin{itemize}
  \item \textbf{Columns:} staff\_id, first\_name, last\_name, email, street, city, postcode, dob, gender, position, department\_name, internal\_phone, room\_number, location, created\_at.
  \item \textbf{Primary key:} staff\_id (AUTO\_INCREMENT).
  \item \textbf{Unique constraints:} email.
  \item \textbf{Foreign keys (outgoing):} none.
  \item \textbf{Indexes:} PRIMARY(staff\_id), UNIQUE(email).
  \item \textbf{Referenced by:} students(advisor\_staff\_id), halls(manager\_staff\_id), inspections(staff\_id).
\end{itemize}

\subsection{\texttt{courses}}
\begin{itemize}
  \item \textbf{Columns:} course\_number, course\_title, instructor\_name, instructor\_phone, instructor\_email, instructor\_room, department\_name, created\_at.
  \item \textbf{Primary key:} course\_number.
  \item \textbf{Foreign keys (outgoing):} none.
  \item \textbf{Indexes:} PRIMARY(course\_number).
  \item \textbf{Referenced by:} students(course\_number).
\end{itemize}

\subsection{\texttt{students}}
\begin{itemize}
  \item \textbf{Columns:} banner\_id, first\_name, last\_name, street, city, postcode, mobile\_phone, email, dob, gender, category, nationality, special\_needs, comments, status, major, minor, course\_number, advisor\_staff\_id, created\_at.
  \item \textbf{Primary key:} banner\_id.
  \item \textbf{Foreign keys (outgoing):}
  \begin{itemize}
    \item course\_number -> courses(course\_number), ON UPDATE CASCADE, ON DELETE SET NULL.
    \item advisor\_staff\_id -> staff(staff\_id), ON UPDATE CASCADE, ON DELETE SET NULL.
  \end{itemize}
  \item \textbf{Indexes:} PRIMARY(banner\_id), idx\_students\_status(status), idx\_students\_category(category), idx\_students\_advisor\_staff\_id(advisor\_staff\_id).
\end{itemize}

\subsection{\texttt{halls}}
\begin{itemize}
  \item \textbf{Columns:} hall\_id, hall\_name, street, city, postcode, telephone, manager\_staff\_id, created\_at.
  \item \textbf{Primary key:} hall\_id (AUTO\_INCREMENT).
  \item \textbf{Foreign keys (outgoing):} manager\_staff\_id -> staff(staff\_id), ON UPDATE CASCADE, ON DELETE RESTRICT.
  \item \textbf{Indexes:} PRIMARY(hall\_id), idx\_halls\_manager\_id(manager\_staff\_id).
\end{itemize}

\subsection{\texttt{apartments}}
\begin{itemize}
  \item \textbf{Columns:} apartment\_id, street, city, postcode, number\_of\_bedrooms, created\_at.
  \item \textbf{Primary key:} apartment\_id (AUTO\_INCREMENT).
  \item \textbf{Checks:} number\_of\_bedrooms IN (3,4,5).
  \item \textbf{Foreign keys (outgoing):} none.
  \item \textbf{Indexes:} PRIMARY(apartment\_id).
\end{itemize}

\subsection{\texttt{rooms}}
\begin{itemize}
  \item \textbf{Columns:} place\_number, room\_number, monthly\_rent, hall\_id, apartment\_id, created\_at.
  \item \textbf{Primary key:} place\_number.
  \item \textbf{Foreign keys (outgoing):}
  \begin{itemize}
    \item hall\_id -> halls(hall\_id), ON UPDATE CASCADE, ON DELETE CASCADE.
    \item apartment\_id -> apartments(apartment\_id), ON UPDATE CASCADE, ON DELETE CASCADE.
  \end{itemize}
  \item \textbf{Checks:} monthly\_rent >= 0.
  \item \textbf{Triggers:} \texttt{trg\_rooms\_before\_insert} and \texttt{trg\_rooms\_before\_update} enforce XOR parent rule (exactly one of hall\_id/apartment\_id).
  \item \textbf{Indexes:} PRIMARY(place\_number), idx\_rooms\_hall\_id(hall\_id), idx\_rooms\_apartment\_id(apartment\_id).
\end{itemize}

\subsection{\texttt{leases}}
\begin{itemize}
  \item \textbf{Columns:} lease\_id, banner\_id, place\_number, duration\_semesters, includes\_summer\_semester, date\_enter, date\_leave, created\_at.
  \item \textbf{Primary key:} lease\_id (AUTO\_INCREMENT).
  \item \textbf{Foreign keys (outgoing):}
  \begin{itemize}
    \item banner\_id -> students(banner\_id), ON UPDATE CASCADE, ON DELETE CASCADE.
    \item place\_number -> rooms(place\_number), ON UPDATE CASCADE, ON DELETE RESTRICT.
  \end{itemize}
  \item \textbf{Indexes:} PRIMARY(lease\_id), idx\_leases\_banner\_id(banner\_id), idx\_leases\_place\_number(place\_number), idx\_leases\_summer(includes\_summer\_semester).
\end{itemize}

\subsection{\texttt{invoices}}
\begin{itemize}
  \item \textbf{Columns:} invoice\_id, lease\_id, semester, amount\_due, due\_date, date\_paid, payment\_method, first\_reminder\_date, second\_reminder\_date, created\_at.
  \item \textbf{Primary key:} invoice\_id (AUTO\_INCREMENT).
  \item \textbf{Foreign keys (outgoing):} lease\_id -> leases(lease\_id), ON UPDATE CASCADE, ON DELETE CASCADE.
  \item \textbf{Checks:} amount\_due >= 0.
  \item \textbf{Indexes:} PRIMARY(invoice\_id), idx\_invoices\_lease\_id(lease\_id), idx\_invoices\_due\_date(due\_date), idx\_invoices\_date\_paid(date\_paid).
\end{itemize}

\subsection{\texttt{next\_of\_kin}}
\begin{itemize}
  \item \textbf{Columns:} kin\_id, banner\_id, name, relationship, street, city, postcode, phone, created\_at.
  \item \textbf{Primary key:} kin\_id (AUTO\_INCREMENT).
  \item \textbf{Unique constraints:} banner\_id (one-to-one-ish from students to kin rows).
  \item \textbf{Foreign keys (outgoing):} banner\_id -> students(banner\_id), ON UPDATE CASCADE, ON DELETE CASCADE.
  \item \textbf{Indexes:} PRIMARY(kin\_id), UNIQUE(banner\_id).
\end{itemize}

\subsection{\texttt{inspections}}
\begin{itemize}
  \item \textbf{Columns:} inspection\_id, apartment\_id, staff\_id, inspection\_date, is\_satisfactory, comments, created\_at.
  \item \textbf{Primary key:} inspection\_id (AUTO\_INCREMENT).
  \item \textbf{Foreign keys (outgoing):}
  \begin{itemize}
    \item apartment\_id -> apartments(apartment\_id), ON UPDATE CASCADE, ON DELETE CASCADE.
    \item staff\_id -> staff(staff\_id), ON UPDATE CASCADE, ON DELETE RESTRICT.
  \end{itemize}
  \item \textbf{Indexes:} PRIMARY(inspection\_id), idx\_inspections\_apartment\_id(apartment\_id), idx\_inspections\_staff\_id(staff\_id), idx\_inspections\_sat(is\_satisfactory).
\end{itemize}

\section{Foreign Key Matrix (All Tables)}
\begin{longtable}{>{\raggedright\arraybackslash}p{0.27\textwidth} >{\raggedright\arraybackslash}p{0.25\textwidth} >{\raggedright\arraybackslash}p{0.12\textwidth} >{\raggedright\arraybackslash}p{0.12\textwidth} >{\raggedright\arraybackslash}p{0.16\textwidth}}
\toprule
\textbf{Child Column} & \textbf{Parent Reference} & \textbf{ON UPDATE} & \textbf{ON DELETE} & \textbf{Purpose} \\
\midrule
\endhead
students.course\_number & courses(course\_number) & CASCADE & SET NULL & Student-course mapping \\
students.advisor\_staff\_id & staff(staff\_id) & CASCADE & SET NULL & Adviser relationship \\
halls.manager\_staff\_id & staff(staff\_id) & CASCADE & RESTRICT & Hall manager assignment \\
rooms.hall\_id & halls(hall\_id) & CASCADE & CASCADE & Room under hall \\
rooms.apartment\_id & apartments(apartment\_id) & CASCADE & CASCADE & Room under apartment \\
leases.banner\_id & students(banner\_id) & CASCADE & CASCADE & Lease owner \\
leases.place\_number & rooms(place\_number) & CASCADE & RESTRICT & Occupancy place reference \\
invoices.lease\_id & leases(lease\_id) & CASCADE & CASCADE & Billing tied to lease \\
next\_of\_kin.banner\_id & students(banner\_id) & CASCADE & CASCADE & Emergency contact link \\
inspections.apartment\_id & apartments(apartment\_id) & CASCADE & CASCADE & Apartment inspection link \\
inspections.staff\_id & staff(staff\_id) & CASCADE & RESTRICT & Inspector accountability \\
\bottomrule
\end{longtable}

\section{Index Catalog and Query Optimization Impact}
\begin{longtable}{>{\raggedright\arraybackslash}p{0.28\textwidth} >{\raggedright\arraybackslash}p{0.18\textwidth} >{\raggedright\arraybackslash}p{0.20\textwidth} >{\raggedright\arraybackslash}p{0.26\textwidth}}
\toprule
\textbf{Index Name} & \textbf{Table} & \textbf{Columns} & \textbf{Typical API/Report Usage} \\
\midrule
\endhead
PRIMARY & staff & staff\_id & Staff CRUD identity lookups \\
UNIQUE(email) & staff & email & Avoid duplicate staff identities \\
PRIMARY & courses & course\_number & Course CRUD identity lookups \\
PRIMARY & students & banner\_id & Student CRUD identity lookups \\
idx\_students\_status & students & status & \texttt{v\_waiting\_list} filter \\
idx\_students\_category & students & category & \texttt{v\_student\_category\_counts} grouping efficiency \\
idx\_students\_advisor\_staff\_id & students & advisor\_staff\_id & Adviser joins in \texttt{v\_student\_advisers} \\
PRIMARY & halls & hall\_id & Hall CRUD identity lookups \\
idx\_halls\_manager\_id & halls & manager\_staff\_id & Hall-manager joins \\
PRIMARY & apartments & apartment\_id & Apartment CRUD identity lookups \\
PRIMARY & rooms & place\_number & Room CRUD identity lookups \\
idx\_rooms\_hall\_id & rooms & hall\_id & Hall room lookups and report generation \\
idx\_rooms\_apartment\_id & rooms & apartment\_id & Apartment-based joins in inspections and invoices views \\
PRIMARY & leases & lease\_id & Lease CRUD identity lookups \\
idx\_leases\_banner\_id & leases & banner\_id & Student lease history, rent-paid aggregates \\
idx\_leases\_place\_number & leases & place\_number & Place occupancy joins \\
idx\_leases\_summer & leases & includes\_summer\_semester & Summer lease filtering \\
PRIMARY & invoices & invoice\_id & Invoice CRUD identity lookups \\
idx\_invoices\_lease\_id & invoices & lease\_id & Lease-invoice join path \\
idx\_invoices\_due\_date & invoices & due\_date & Due-date filtering in unpaid invoices report \\
idx\_invoices\_date\_paid & invoices & date\_paid & Paid/unpaid partition scans \\
PRIMARY & next\_of\_kin & kin\_id & Next-of-kin CRUD identity lookups \\
UNIQUE(banner\_id) & next\_of\_kin & banner\_id & Enforces max one kin row per student \\
PRIMARY & inspections & inspection\_id & Inspection CRUD identity lookups \\
idx\_inspections\_apartment\_id & inspections & apartment\_id & Apartment inspection history joins \\
idx\_inspections\_staff\_id & inspections & staff\_id & Inspector-based filtering \\
idx\_inspections\_sat & inspections & is\_satisfactory & Unsatisfactory inspection report filter \\
\bottomrule
\end{longtable}

\section{Deployment and Runtime Notes}
\subsection{Local Execution}
\begin{enumerate}
  \item Apply schema to MySQL and optional seed.
  \item Configure backend \texttt{backend/.env}.
  \item Run backend with Uvicorn and frontend with Vite.
\end{enumerate}

\subsection{Cloud Split Deployment}
\begin{itemize}
  \item Frontend: Vercel or Netlify
  \item Backend: Render or Railway
  \item Database: MySQL provider (Aiven / Railway / other managed MySQL)
\end{itemize}

\subsection{TLS and Managed DB Logs}
Startup messages such as "ready for connections" and "TLS enabled" indicate healthy DB startup. Self-signed CA warning is common for managed instances and handled by TLS mode/cert configuration.

\section{Testing and Quality Assurance}
\subsection{Route-Level Tests}
\begin{lstlisting}[language={}]
source .venv/bin/activate
PYTHONPATH=backend python -m pytest backend/tests -q
\end{lstlisting}

\subsection{Runtime Smoke Test}
\begin{lstlisting}[language={}]
source .venv/bin/activate
python backend/tests/runtime_smoke.py
\end{lstlisting}

\subsection{What Is Validated}
\begin{itemize}
  \item endpoint availability
  \item auth behavior and token flow
  \item CRUD/report response correctness
\end{itemize}

\section{Known Constraints and Future Improvements}
\begin{itemize}
  \item Current CORS policy is open; production should restrict origins.
  \item Demo users are in-memory seeded; production should store users in DB.
  \item Add stored procedures/functions if rubric explicitly requires them.
  \item Add advanced auditing, pagination, and rate limiting.
\end{itemize}

\section{Final Conclusion}
This project is a complete DBMS-focused software system that demonstrates relational modeling, integrity controls, transactional correctness, secure service APIs, and practical frontend integration.

From first user action to committed row, the flow is fully traceable:
\begin{enumerate}
  \item typed frontend interaction
  \item authenticated API call
  \item validated backend logic
  \item transactional SQL execution
  \item consistent and durable relational state
\end{enumerate}

It therefore satisfies both engineering and academic DBMS expectations: normalized schema design, ACID-aware operations, SQL reporting via views, and end-to-end application delivery.

\appendix
\section{Full Report Endpoint Quick List}
\begin{lstlisting}[language={}]
GET /api/reports/hall-managers
GET /api/reports/student-leases
GET /api/reports/summer-leases
GET /api/reports/student-rent-paid/{banner_id}
GET /api/reports/unpaid-invoices?due_before=YYYY-MM-DD
GET /api/reports/unsatisfactory-inspections
GET /api/reports/hall-student-rooms/{hall_id}
GET /api/reports/waiting-list
GET /api/reports/student-category-counts
GET /api/reports/students-without-kin
GET /api/reports/student-adviser/{banner_id}
GET /api/reports/rent-stats
GET /api/reports/hall-place-counts
GET /api/reports/senior-staff
\end{lstlisting}

\section{Generic CRUD Pattern Quick List}
\begin{lstlisting}[language={}]
GET    /api/{entity}
GET    /api/{entity}/{record_id}
POST   /api/{entity}
PUT    /api/{entity}/{record_id}
DELETE /api/{entity}/{record_id}
\end{lstlisting}

\end{document}
